%% =============================================================================
%% sm_b1_scores.tex — SM-B, Subsection B.1: Score Functions
%% Body content only (\subsection level); parent structure in sm_b.tex
%% Primary source: dev/log/03_research-notes-deep-mathematical-properties.qmd
%%   lines 620--689 (extensive/intensive scores, Fisher information)
%% Corrections applied: T1-2 (truncated mu-score sign: + -> -)
%% =============================================================================

\subsection{Score functions}
\label{smb:scores}

We derive the score functions for each component of the hurdle
beta-binomial log-likelihood.  These expressions serve three roles:
\textup{(i)}~verification of the automatic differentiation gradients
used in Stan~\citep{CarpenterEtAl2017}, \textup{(ii)}~diagnostics via
score magnitude monitoring,
and \textup{(iii)}~the theoretical basis for the sandwich variance
estimator~\citep{Binder1983} in~\cref{app:survey}.
Recall the parameterization $a = \mu\kappa$, $b = (1-\mu)\kappa$, the
zero probability $p_0$ from~\eqref{eq:p0}, and the auxiliary quantities
$\Lambda$, $\Lambda_2$ defined in~\cref{sma:eq-Lambda}.

%% ---------------------------------------------------------------------------
%%  B.1.1  Extensive-margin score
%% ---------------------------------------------------------------------------

\begin{proposition}[Extensive-margin score]
\label{smb:prop-score-ext}
  The score of the extensive-margin log-likelihood
  $\ell_i^{(1)} = z_i\log q_i + (1-z_i)\log(1-q_i)$ with respect to
  the linear predictor
  $\eta_i^{(1)} = \bx_i^\t\balpha + \bx_i^{(r)\t}\bdelta_{1,s}$ is
  \begin{equation}\label{smb:eq-score-ext}
    \frac{\partial\ell_i^{(1)}}{\partial\eta_i^{(1)}}
    = z_i - q_i,
  \end{equation}
  where $q_i = \expit(\eta_i^{(1)})$.  The Hessian is non-random:
  \begin{equation}\label{smb:eq-hess-ext}
    \frac{\partial^2\ell_i^{(1)}}{\partial(\eta_i^{(1)})^2}
    = -\,q_i(1-q_i).
  \end{equation}
\end{proposition}

\begin{proof}
  Standard logistic regression calculus.  Since
  $\partial q_i/\partial\eta_i^{(1)} = q_i(1-q_i)$,
  differentiating $\ell_i^{(1)}$ yields
  $z_i/q_i \cdot q_i(1-q_i) - (1-z_i)/(1-q_i) \cdot q_i(1-q_i)
  = z_i(1-q_i) - (1-z_i)q_i = z_i - q_i$.
  The second derivative is $-q_i(1-q_i)$, which depends only on
  parameters and is therefore non-random.
\end{proof}

\begin{remark}\label{smb:rem-ext-info}
  The non-randomness of~\eqref{smb:eq-hess-ext} implies that the
  observed and expected Fisher information coincide for the extensive
  margin:
  $\cl{I}_{\textup{ext},i} = q_i(1-q_i)$.
  This simplifies the sandwich estimator, since the Hessian bread matrix
  for Part~1 contains no stochastic terms.
\end{remark}

%% ---------------------------------------------------------------------------
%%  B.1.2  Intensive-margin score for mu
%% ---------------------------------------------------------------------------

\begin{proposition}[Intensive-margin score for $\mu$]
\label{smb:prop-score-mu}
  Define the (untruncated) beta-binomial score function
  \citep[following the parameterization of][]{Prentice1986}
  \begin{equation}\label{smb:eq-Sbb}
    S_{\textup{BB}}(y;\, \mu,\kappa,n)
    = \kappa\bigl[\psi(y+a) - \psi(n-y+b) - \psi(a) + \psi(b)\bigr],
  \end{equation}
  where $\psi(\cdot)$ is the digamma function and $a = \mu\kappa$,
  $b = (1-\mu)\kappa$.  The score of the zero-truncated
  log-likelihood\/ $\ell_i^{(2+)} = \log f_{\textup{BB}}(y_i)
  - \log(1-p_{0,i})$ with respect to $\mu$ is
  \begin{equation}\label{smb:eq-score-mu}
    \boxed{%
      \frac{\partial\ell_i^{(2+)}}{\partial\mu}
      = S_{\textup{BB}}(y_i)
        \;-\; \frac{p_{0,i}\,\Lambda_i}{1-p_{0,i}}
    }.
  \end{equation}
\end{proposition}

\begin{proof}
  The decomposition $\ell_i^{(2+)} = \log f_{\textup{BB}}(y_i)
  - \log(1-p_{0,i})$ gives two terms.

  \medskip\noindent\textit{Term~1 (BB score).}\enspace
  Using $\partial a/\partial\mu = \kappa$ and
  $\partial b/\partial\mu = -\kappa$, differentiation of the
  beta-binomial log-PMF yields
  \[
    \frac{\partial\log f_{\textup{BB}}}{\partial\mu}
    = \kappa\bigl[\psi(y+a) - \psi(n-y+b) - \psi(a) + \psi(b)\bigr]
    = S_{\textup{BB}}(y).
  \]

  \medskip\noindent\textit{Term~2 (truncation correction).}\enspace
  Since $\ln p_0 = \sum_{j=0}^{n-1}\ln[(1-\mu)\kappa + j]
  - \ln[\kappa + j]$, we have $\partial p_0/\partial\mu = -p_0\Lambda$
  (see~\cref{sma:eq-Lambda}), and therefore
  \[
    \frac{\partial\log(1-p_0)}{\partial\mu}
    = \frac{p_0\,\Lambda}{1-p_0}
    > 0.
  \]
  Combining, $\partial\ell^{(2+)}/\partial\mu
  = S_{\textup{BB}}(y) - p_0\Lambda/(1-p_0)$.
\end{proof}

%% ---------------------------------------------------------------------------
%%  B.1.3  Score for log kappa
%% ---------------------------------------------------------------------------

\begin{proposition}[Score for $\log\kappa$]
\label{smb:prop-score-kappa}
  The score of the zero-truncated log-likelihood with respect to the
  log-dispersion parameter is
  \begin{equation}\label{smb:eq-score-kappa}
    \boxed{%
    \begin{aligned}
      \frac{\partial\ell_i^{(2+)}}{\partial\log\kappa}
      &= \kappa\Bigl[
           \mu\bigl[\psi(y_i\!+\!a)-\psi(a)\bigr]
           + (1\!-\!\mu)\bigl[\psi(n_i\!-\!y_i\!+\!b)-\psi(b)\bigr]
      \\
      &\qquad\quad
           + \psi(\kappa)-\psi(n_i\!+\!\kappa)
         \Bigr]
       - \frac{p_{0,i}\,\mu\kappa}{1-p_{0,i}}
         \sum_{j=1}^{n_i-1}
         \frac{j}{(b+j)(\kappa+j)}.
    \end{aligned}
    }
  \end{equation}
\end{proposition}

\begin{proof}
  With $\phi = \log\kappa$ and the chain rule
  $\partial/\partial\phi = \kappa\,\partial/\partial\kappa$, the
  untruncated BB log-PMF gives the first line
  via $\partial a/\partial\kappa = \mu$ and
  $\partial b/\partial\kappa = 1-\mu$.  The truncation correction
  requires $\partial\log(1-p_0)/\partial\log\kappa$.  Writing
  $\ln p_0 = \sum_{j=0}^{n-1}\ln(b+j) - \ln(\kappa+j)$ and
  differentiating with respect to $\kappa$ (noting
  $\partial b/\partial\kappa = 1-\mu$) yields the second line after
  rearranging the telescoping sums.
\end{proof}

%% ---------------------------------------------------------------------------
%%  B.1.4  Expected Fisher information for mu
%% ---------------------------------------------------------------------------

\begin{proposition}[Expected information for $\mu$]
\label{smb:prop-info-mu}
  The negative expected second derivative of the zero-truncated
  log-likelihood with respect to $\mu$ is
  \begin{multline}\label{smb:eq-info-mu}
    -\E_{Y|Y>0}\!\left[
      \frac{\partial^2\ell^{(2+)}}{\partial\mu^2}
    \right]
    = \frac{\kappa^2}{1-p_0}\,
      \E_{Y\sim\BetaBin}\bigl[\psi_1(Y+a) + \psi_1(n\!-\!Y+b)\bigr] \\
      - \kappa^2\bigl[\psi_1(a)+\psi_1(b)\bigr]
      + T'(\mu),
  \end{multline}
  where $\psi_1(\cdot) = d\psi/dx$ is the trigamma function and
  \begin{equation}\label{smb:eq-Tprime}
    T'(\mu)
    = \frac{p_0\,\Lambda_2}{1-p_0}
      - \frac{p_0\,\Lambda^2}{(1-p_0)^2},
    \qquad
    \Lambda_2
    = \kappa^2 \sum_{j=0}^{n-1}\frac{1}{(b+j)^2}.
  \end{equation}
\end{proposition}

\begin{proof}
  Differentiating the score~\eqref{smb:eq-score-mu} with respect to
  $\mu$ produces two contributions.

  \medskip\noindent\textit{The $S_{\textup{BB}}$ term.}\enspace
  Direct differentiation gives
  $\partial S_{\textup{BB}}/\partial\mu
  = \kappa^2[\psi_1(y+a) + \psi_1(n-y+b) - \psi_1(a) - \psi_1(b)]$.

  \medskip\noindent\textit{The truncation term.}\enspace
  Set $T(\mu) = p_0\Lambda/(1-p_0)$.  The quotient rule with
  $\partial p_0/\partial\mu = -p_0\Lambda$ gives
  \begin{align*}
    T'(\mu)
    &= \frac{p_0(\Lambda_2 - \Lambda^2)(1-p_0)
             - p_0^2\Lambda^2}{(1-p_0)^2}
    = \frac{p_0\Lambda_2(1-p_0) - p_0\Lambda^2}{(1-p_0)^2} \\
    &= \frac{p_0\,\Lambda_2}{1-p_0}
       - \frac{p_0\,\Lambda^2}{(1-p_0)^2}.
  \end{align*}

  \medskip\noindent\textit{Expectation under truncation.}\enspace
  Taking expectations under $Y \mid Y > 0$ and converting via the
  identity $\E_{Y|Y>0}[g(Y)] = \E_{Y\sim\BetaBin}[g(Y)]/(1-p_0)$
  yields~\eqref{smb:eq-info-mu}.  The non-stochastic terms
  $-\kappa^2[\psi_1(a)+\psi_1(b)]$ and $T'(\mu)$ pass through the
  expectation unchanged.
\end{proof}

%% ---------------------------------------------------------------------------
%%  B.1.5  Cross-information and information matrix structure
%% ---------------------------------------------------------------------------

\begin{remark}[Cross-information]\label{smb:rem-cross-info}
  Both $\mu$ and $\kappa$ enter $\ell_i^{(2+)}$, so the
  cross-information $\cl{I}_{\mu,\kappa}$ is nonzero.
  The intensive-margin information block is
  \begin{equation}\label{smb:eq-info-block}
    \cl{I}_{\textup{int}+\kappa}
    = \begin{pmatrix}
        \cl{I}_{\textup{int}}   & \cl{I}_{\textup{int},\kappa} \\
        \cl{I}_{\kappa,\textup{int}} & \cl{I}_{\kappa}
      \end{pmatrix},
  \end{equation}
  which is not further block-diagonal.  The intensive-margin regression
  parameters and $\kappa$ must be estimated jointly, but they are fully
  separated from the extensive-margin parameters by the hurdle structure
  (\cref{sec:model}).
\end{remark}

